\chapter{스키마를 계약으로 두는 설계}
\label{ch:contract}

이 장은 proto 메시지 → codec → ClickHouse DDL로 이어지는 3단계를 "같은 계약을
다른 형태로 반복하는 사슬"로 보는 관점을 다룬다. \texttt{venue} 필드 하나를 추가하는
결정이 왜 병렬 테이블이나 별도 마이그레이션 없이 끝나는지가 핵심이며, 이 계약의
견고함이 \Cref{ch:delivery}에서 다루는 정렬 키(sort key) 기반 중복 제거의 전제가 된다.

\todo{§3(스키마 3단계, venue 차원 도입 근거, 정렬 키에 venue를 선두에 두는 이유)를
바탕으로 "계약이 하나여야 하는 이유"를 설명한다.}
